\documentclass[10pt,a4paper]{article}
\usepackage{auditstyle}
\graphicspath{{evidence/}}
\title{Citation screenshot dossier\\[3pt]
\large Exact source pages used in the 2026 audit of \texttt{glimm\_jaffe/m6/main.tex}}
\author{}
\date{9 August 2026}


% --- Page images that are not distributed with this repository ------------------
% The citation page images are third-party copyrighted material and are kept out of
% the public repository (see the README).  Each image is therefore looked up before
% it is included, and a framed placeholder naming the file is drawn when it is
% absent, so this document builds without errors either way.
\newcommand{\missingshot}[1]{%
  \fbox{\begin{minipage}[c][0.28\textheight][c]{0.86\textwidth}%
    \centering
    \textbf{Page image not distributed with this repository}\par\medskip
    \texttt{\detokenize{#1}}\par\medskip
    {\small Third-party copyrighted material.  Regenerate the corpus from your own
     copies of the sources with
     \texttt{fresh\_audit\_2026\_08\_16/render\_evidence.sh}.}%
  \end{minipage}}}
\newcommand{\evidence}[5]{%
  \clearpage
  \phantomsection\addcontentsline{toc}{section}{#1 --- #2}
  {\Large\bfseries #1 --- #2\par}
  \vspace{3pt}
  \textbf{Claim audited.} #3\par
  \textbf{Finding.} #4\par
  \vspace{3pt}
  \begin{center}
  \IfFileExists{evidence/#5}{%
    \includegraphics[width=\textwidth,height=.76\textheight,keepaspectratio]{#5}}{%
    \missingshot{evidence/#5}}
  \end{center}
}

\begin{document}
\maketitle

\begin{resultbox}
This is the page-image evidence companion to the critical audit and the two corrected
Lamport proofs.  Every bibliography item used by M6 is represented.  Load-bearing
citations are shown at the theorem page actually used; contextual citations are explicitly
labelled contextual.  A page image verifies only what appears on that page, not any stronger
claim.

Page numbers called ``PDF page'' count from the local file's first page; ``printed page'' is
the journal/book pagination visible in the image.  All source PDFs and SHA-256 values are
listed in the critical audit report.
\end{resultbox}

\tableofcontents

\section{Evidence index}

\begin{longtable}{P{1.0cm}P{3.2cm}P{5.1cm}P{4.7cm}}
\toprule
ID & Source & Result/claim & Audit use\\
\midrule
\endhead
E01--02 & Lamport & hierarchical numbering and proof layout & method, not mathematical input\\
E03 & GJ I & exact finite propagation/cutoff locality & load-bearing construction input\\
E04--06 & GJ II & ground-state existence/uniqueness; cutoff covariance/local algebra & load-bearing construction inputs\\
E07--08 & GJ III & averaged limit point and local Fock normality & historical claim and canonical-vacuum corollary\\
E09 & GJ IV & Wightman construction statement & contextual construction history\\
E10--11 & GJ 1971 & lower bounds and positivity estimates & corroborating cutoff construction\\
E12--13 & Rosen & semigroup limit and self-adjoint generator & corroborating cutoff construction\\
E14--19 & GJ Expositions & unproved vacuum limit, propagation, domain smoothing & historical and load-bearing; one citation repaired\\
E20--26 & GJS 1974 & cutoff class, scaling, norm, CE2, CE1, uniformity & principal weak-coupling inputs\\
E54 & GJS 1974 & proof of Thm.~1.1.7 from Thms.~1.1.8/1.1.11 (p.~597) & scope of $U_{\rm GJS}$; added 5 Sept.\ 2026\\
E27--32 & GRS 1975 & FKN and Problems 1--2, including note added & load-bearing FKN and historical scope\\
E33--43 & Simon & hypercontractivity, FKN, stability, miscited III.12, interacting FKN & pinpoint correction of three citations\\
E44--45 & AHZ & Euclidean uniqueness/global Markov property & contextual only\\
E46 & BDW & stochastic-quantization invariant-measure uniqueness & contextual; not Hamiltonian uniqueness\\
E47 & RZT & finite-volume scope & contextual scope check\\
E48--50 & GJS 1975 & phase transition and pure phases & contextual; insufficient for omitted M3 premise\\
E51--53 & GJS 1976 & two pure phases and mean-field theorems & contextual; insufficient for omitted M3 premise\\
\bottomrule
\end{longtable}

\evidence{E01}{Lamport worked structured proof, first page}
{The requested notation is a hierarchical proof whose labels expose logical structure.}
{Method evidence.  The companion proofs use labels of the form $\langle1\rangle1$,
$\langle2\rangle1$, with explicit ``Because'' clauses.}
{lamport-example-10.png}

\evidence{E02}{Lamport numbering discussion}
{Hierarchical step numbering is the relevant Lamport convention.}
{Method evidence.  This page supports the notation choice; it is not an input to the
field-theory theorem.}
{lamport-numbering-16.png}

\evidence{E03}{Glimm--Jaffe I, PDF page 6}
{Cutoff dynamics has finite propagation and a local observable depends on the cutoff only
inside the causal region.}
{Load-bearing pass.  The page contains the Trotter/locality argument and the cutoff-dependence
statement used for patching.}
{gj1-locality-06.png}

\evidence{E04}{Glimm--Jaffe II, Theorem 2.2.1}
{The spatially cutoff Hamiltonian has a vacuum vector.}
{Load-bearing pass.  This is the exact existence theorem page.}
{gj2-ground-existence-08.png}

\evidence{E05}{Glimm--Jaffe II, Theorem 2.3.1}
{The cutoff vacuum vector is unique.}
{Load-bearing pass.  Simplicity is used in the slab limit and parity argument.}
{gjii-ground-12.png}

\evidence{E06}{Glimm--Jaffe II, printed page 398}
{The field is independent of the spatial cutoff inside the causal region and transforms
covariantly under translations.}
{Load-bearing pass.  This page states the cutoff-independence condition explicitly.}
{gj2-cutoff-independence-38.png}

\evidence{E07}{Glimm--Jaffe III, printed page 210}
{The physical vacuum construction begins with averaged cutoff states and a weak-star limit
point.}
{Historical pass.  The displayed averaging formula and Theorem 2.1 support the manuscript's
description.}
{gjiii-08.png}

\evidence{E08}{Glimm--Jaffe III, printed page 211}
{The limit-point representation is locally Fock/normal.}
{Historical pass.  Theorems 2.2--2.3 are visible and match the cited conclusion.}
{gjiii-09.png}

\evidence{E09}{Glimm--Jaffe IV, first page}
{The I--IV series establishes limiting Wightman functions and field operators.}
{Contextual pass.  Theorems A and B support the construction-history sentence; no later
proof step depends on this page.}
{gj4-01.png}

\evidence{E10}{Glimm--Jaffe 1971, Theorem 1}
{Cutoff Hamiltonians possess uniform lower bounds in the auxiliary cutoffs.}
{Corroborating pass.  This is not, by itself, the source for every simplicity/domain claim.}
{gj71-lower-bound-02.png}

\evidence{E11}{Glimm--Jaffe 1971, Theorem 2}
{The cited positivity/self-adjointness paper supplies the stated model estimates.}
{Corroborating pass.  M6 should continue to rely on the more precise GJ II/GJS statements
for ground-state simplicity.}
{gj71-theorem2-04.png}

\evidence{E12}{Rosen 1970, semigroup construction}
{The cutoff-free/self-adjoint construction is supported by a convergent semigroup.}
{Corroborating pass.  Theorem 5.3 begins on this page.}
{rosen-semigroup-22.png}

\evidence{E13}{Rosen 1970, self-adjoint generator/core}
{The limiting semigroup generator is positive self-adjoint and agrees with the Hamiltonian
on the specified core.}
{Corroborating pass.  The visible continuation supplies the precise operator conclusion.}
{rosen-selfadjoint-23.png}

\evidence{E14}{Glimm--Jaffe Expositions, PDF page 54}
{The original construction did not prove convergence of cutoff vacuum expectation values.}
{Historical pass.  The relevant discussion begins on this page.}
{gjexp-054.png}

\evidence{E15}{Glimm--Jaffe Expositions, PDF page 55}
{The limit of vacuum expectations ``has not been proved'' and locality uses finite
propagation.}
{Historical/load-bearing pass.  The exact limiting-status language and surrounding
propagation argument are visible.}
{gjexp-055.png}

\evidence{E16}{Glimm--Jaffe Expositions, PDF page 56}
{Cutoff dynamics is patched into a local cutoff-independent dynamics.}
{Load-bearing pass.  This is the continuation of the construction shown in E15.}
{gjexp-056.png}

\evidence{E17}{Glimm--Jaffe Expositions, Theorem 7.3}
{M6 cites Theorem 7.3 for semigroup domain smoothing.}
{Citation mismatch.  The visible statement assumes the multiplication potential $V$ is
bounded below, a hypothesis M6 does not verify for its Wick interaction.}
{gjexp-domain-165.png}

\evidence{E18}{Glimm--Jaffe Expositions, Lemma 7.4}
{The proof mechanism for domain smoothing uses convergence of truncated multiplication
operators times the semigroup.}
{Supporting page.  The lemma belongs to the general Theorem 7.3 argument; the model-specific
result is on the next page.}
{gjexp-domain-166.png}

\evidence{E19}{Glimm--Jaffe Expositions, Theorem 7.4}
{For the $P(\varphi)_2$ Hamiltonian, $\Ran e^{-tH(g)}$ lies in the domain of every
relevant $V'$.}
{Load-bearing repair.  Cite this theorem directly for M6 line 214.}
{gjexp-domain-167.png}

\evidence{E20}{GJS 1974, printed page 588}
{The spatially cutoff Hamiltonian is essentially self-adjoint and has a simple ground
eigenvalue.}
{Load-bearing pass.  The source states the cutoff model and simplicity explicitly.}
{gjs-source-05.png}

\evidence{E21}{GJS 1974, printed page 589}
{The Euclidean cutoff class permits compactly supported $0\leq h\leq1$ and defines
$dq_h$.}
{Load-bearing pass.  The slab cutoff $\mathbf1_{[-T,T]}\otimes g$ lies in this measurable
class.}
{gjs-source-06.png}

\evidence{E22}{GJS 1974, printed page 590}
{The dimensionless coupling is $\lambda/m_0^2$, by an exact unitary dilation.}
{Load-bearing pass.  This page supplies the scaling omitted from M6's proof.}
{gjs-scaling-07.png}

\evidence{E23}{GJS 1974, printed page 593}
{The observable class is the product of individually Wick-ordered powers, with bounded
per-point degree.}
{Load-bearing pass.  This validates the $V(h)^2$ observable without Wick recombination.}
{gjs-theorems-10.png}

\evidence{E24}{GJS 1974, printed page 594}
{The GJS norms have one overall $K_1$, local weights $K_2^{2N(\Delta)}N(\Delta)!$, and
Theorem 1.1.7 has constants independent of $h,Q_1,Q_2$.}
{Load-bearing pass.  This closes the norm-prefactor and CE2 gates.}
{gjs-theorems-11.png}

\evidence{E25}{GJS 1974, printed page 595}
{Theorem 1.1.8 bounds sharp-time monomials uniformly as the cutoff tends to one.  The same
page, (1.1.12), applies Theorem 1.1.7 at every cutoff $h_1+\alpha(h_2-h_1)$ of an affine
path between arbitrary cutoff functions $h_1,h_2$ (``We apply Theorem 1.1.7 to each term
in the sum'').}
{Load-bearing pass with E26 and E54.  E26 supplies the explicit space-cutoff uniformity
statement; E54 shows the printed proof of Theorem 1.1.7 consuming Theorem 1.1.8 at the
same $h$; the affine-path application fixes the scope at which the source uses the
quantifier (Part~I, Assumption A8(vii$'$), added 5 September 2026).}
{gjs-theorems-12.png}

\evidence{E26}{GJS 1974, printed page 629}
{The estimates of Theorem 1.1.8 are uniform in the space cutoff $h$.}
{Load-bearing pass.  Together with E24's direct independence statement for CE2, the
cutoff-uniform source gate closes.}
{gjs-uniform-46.png}

\evidence{E27}{GRS 1975, Theorem II.14}
{A cutoff Feynman--Kac--Nelson formula identifies Euclidean and Hamiltonian matrix
elements.}
{Load-bearing pass.  This is an appropriate pinpoint replacement for the broad GRS
citation in M6.}
{grs-fkn-39.png}

\evidence{E28}{GRS 1975, Theorem II.16}
{The general interacting FKN formula used for slab transfer is available.}
{Load-bearing pass.  The theorem begins on this page.}
{grs-fkn-40.png}

\evidence{E29}{GRS 1975, printed page 125}
{Problem 1 asks for local $L^p$ estimates on the vacuum density.}
{Historical pass.  This is Euclidean and is not proved by M6.}
{grs-16.png}

\evidence{E30}{GRS 1975, printed page 126}
{Problem 2 asks for convergence of finite-volume Euclidean states on bounded regions.}
{Historical pass with scope qualification: M6 proves a Hamiltonian local-algebra analogue,
not literally this Euclidean problem.}
{grs-17.png}

\evidence{E31}{GRS 1975, printed page 127}
{The surrounding open-problem list fixes the historical context.}
{Contextual evidence.  No load-bearing theorem is imported from this page.}
{grs-18.png}

\evidence{E32}{GRS 1975, printed page 128, note added}
{Later developments affected the status of the first problems.}
{Historical qualification.  The note records Newman's small-coupling $p=1$ result, so
unqualified claims that every form remained open are too broad.}
{grs-19.png}

\evidence{E33}{Simon, Theorem I.17}
{Nelson hypercontractivity supplies the exact $L^p\to L^q$ condition used in path
regularity.}
{Load-bearing repair.  This is the direct theorem M6 should cite at line 542.}
{simon-hyper-056.png}

\evidence{E34}{Simon, Theorem III.6}
{The free-field FKN formula relates time slices and the free Hamiltonian semigroup.}
{Load-bearing repair.  This is the correct free FKN theorem.}
{simon-free-fkn-111.png}

\evidence{E35}{Simon, Theorem III.12}
{M6 cites this theorem as FKN.}
{Citation failure.  The page shows a product-of-projections/Sobolev-region estimate, not
an FKN identity.}
{simon-wrong-iii12-119.png}

\evidence{E36}{Simon, Theorem III.17, first page}
{The time-region form of hypercontractivity supplies the transfer estimate.}
{Load-bearing repair.  The theorem statement begins here.}
{simon-timeslice-125.png}

\evidence{E37}{Simon, Theorem III.17, continuation}
{The hypercontractive exponent condition and proof continue.}
{Load-bearing supporting page for E36.}
{simon-timeslice-126.png}

\evidence{E38}{Simon, Theorem V.7}
{Finite-volume exponential stability/integrability of the interaction weights is available.}
{Load-bearing pass for the actual role at M6 line 557.}
{simon-exp-174.png}

\evidence{E39}{Simon, Theorem V.10, first page}
{M6 cites V.10 as if it were the direct free hypercontractivity theorem.}
{Pinpoint mismatch.  This result is a later transfer-matrix estimate that uses earlier
hypercontractivity; cite I.17/III.17 for the displayed inequality.}
{simon-transfer-178.png}

\evidence{E40}{Simon, Theorem V.10, continuation}
{The transfer estimate's actual scope.}
{Supporting mismatch evidence for E39.}
{simon-transfer-179.png}

\evidence{E41}{Simon, Theorem V.12}
{M6 cites V.12 both for hypercontractivity and for a common-core/self-adjointness role.}
{Mixed verdict.  The hypercontractivity citation is wrong; the essential-self-adjointness/
core citation near M6 line 593 is consistent with this page.}
{simon-transfer-180.png}

\evidence{E42}{Simon, Corollary V.14}
{The interacting semigroup admits the FKN representation.}
{Load-bearing repair for M6 lines 305 and 536.}
{simon-fkn-184.png}

\evidence{E43}{Simon, Theorem V.15}
{The full interacting FKN formula is stated.}
{Load-bearing repair.  This is the strongest direct Simon page matching the formula used.}
{simon-fkn-185.png}

\evidence{E44}{AHZ 1989, abstract}
{Euclidean uniqueness and global Markov results exist for general polynomial interactions
under specified conditions.}
{Contextual pass.  This does not itself prove Hamiltonian cutoff-net convergence.}
{ahz-01.png}

\evidence{E45}{AHZ 1989, uniqueness conclusion}
{The paper's Gibbs-state uniqueness conclusion and transition to the global Markov property.}
{Contextual pass.  The page fixes the actual Euclidean scope of the citation.}
{ahz-uniqueness-24.png}

\evidence{E46}{Bauerschmidt--Dagallier--Weber, first page}
{The modern uniqueness citation concerns $\Phi^4_2$ stochastic quantization.}
{Contextual scope finding.  The abstract proves uniqueness of the invariant measure above
critical temperature, not uniqueness of Hamiltonian ground states or convergence of M6's
cutoff net.}
{bdw-01.png}

\evidence{E47}{Rodriguez Zarate--Thiemann, first page}
{The closest cited Hamiltonian-renormalisation work is finite volume.}
{Contextual pass.  The abstract explicitly states the finite-volume scope.}
{rzt-01.png}

\evidence{E48}{GJS 1975, first page}
{The paper proves a phase transition for specified two-dimensional quantum fields.}
{Contextual pass.  It supports the broad phase-transition statement, not the full omitted
M3 phase-decomposition premise.}
{gjs75-01.png}

\evidence{E49}{GJS 1975, second page}
{The discussion identifies two distinct pure phases and clustering within them.}
{Contextual partial support for M3.  The stronger classification of every parity-invariant
Hamiltonian ground state is not supplied merely by this prose.}
{gjs75-phases-02.png}

\evidence{E50}{GJS 1975, Theorem 1.1 page}
{The paper's precise phase-transition theorem has model-specific hypotheses and an order
parameter conclusion.}
{Scope check.  M6 must quote the exact theorem if it wishes to import it; a general
``coexistence'' phrase cannot replace M3's hypotheses.}
{gjs75-theorem-03.png}

\evidence{E51}{GJS 1976, first page}
{The mean-field expansion studies a selected pure phase and ground state.}
{Contextual pass.  This supports detailed strong-coupling phase analysis.}
{gjs76-01.png}

\evidence{E52}{GJS 1976, second page}
{The paper states mass gaps in each of two pure phases.}
{Important scope distinction.  Gaps in infinite-volume pure-phase representations are not
the same proposition as a uniform lower bound for the finite-cutoff tunnelling gap.}
{gjs76-phases-02.png}

\evidence{E53}{GJS 1976, Theorem 2.1 page}
{The detailed mean-field theorem begins under explicit hypotheses.}
{Contextual scope check.  Nothing on this page removes the need to state M3 Hypothesis 4.1
when invoking M3 Theorem 4.2.}
{gjs76-theorem-05.png}

\evidence{E54}{GJS 1974, printed page 597}
{``Proof of Theorem 1.1.7, assuming Theorems 1.1.8 and 1.1.11'' (begins at the foot of
p.~596): $Q=\int A\,dq_h+O(1)\|A\|e^{-m_0(1-\varepsilon)T}$, then ``the proof follows from
Theorem 1.1.8, (1.1.16) and the above determination of $Q$''.}
{Load-bearing pass for the scope of $U_{\rm GJS}$ (added 5 September 2026).  The
$h$-independent constant of Theorem 1.1.7 is obtained from Theorem 1.1.8 at the same $h$;
accepting Theorem 1.1.7 as printed at non-rectangular slab cutoffs therefore commits to
the $h$-uniform constant of Theorem 1.1.8 on the same class.  Together with E25's
affine-path application and E26, this fixes $U_{\rm GJS}$ as the quantifier the source
states and uses (Part~I, Assumption A8(vi), (vii$'$) and Scope paragraph).}
{gjs-proof-14.png}

\clearpage
\section{Dossier conclusion}

The page-image audit closes the two principal source gates: GJS CE1/CE2 have the claimed
norm structure and cutoff independence, and appropriate FKN/hypercontractivity/domain
theorems exist.  It also proves that the theorem numbers presently attached to those three
imports in M6 are inaccurate.  Finally, it makes the subsidiary-corollary defect visible:
the phase-transition sources establish substantial strong-coupling results, but M6 does not
state the exact phase-space classification assumed by its internal dichotomy theorem.

\end{document}
